% GENERATED FILE. Edit canonical lessons, then run npm run generate:books.
% canonical-source-hash: fnv1a64:ceaf475ef72bb02f
% canonical-lessons: SA-C213-sadyah, SA-C213-tatksanam, SA-C213-prak, SA-C213-anantaram, SA-C213-yavat

\chapter{At once, Instantly, Formerly, Afterwards, As long as}
\label{ch:sa-at-once-instantly-formerly-afterwards-as-long-as}
\hlchaptermodality{213}

\begin{chapteropening}
\textbf{By the end of this chapter:} \emph{I can say at once, instantly, formerly, afterwards and as long as.}

Complete the last lesson of chapter 213: I can say at once, instantly, formerly, afterwards and as long as.
\end{chapteropening}

\section[\texorpdfstring{\sk{सद्यः} (sadyaḥ)}{सद्यः (sadyaḥ)}]{\sk{सद्यः} (sadyaḥ) — at once}

\label{lesson:SA-C213-sadyah}



\begin{quote}
Before the new one: say the Sanskrit for a dream, then the Sanskrit for satisfaction.
\end{quote}

\subsection*{You'll want to know: \sk{सद्यः}}
\textbf{\sk{सद्यः}} — \emph{sadyaḥ} — \textquotedblleft{}at once\textquotedblright{}.

\begin{grammarlens}[title={What you've built}]
One more word for this chapter.
\end{grammarlens}

\subsection*{Your turn}
\begin{itemize}
\raggedright
  \item \emph{Say it:} \emph{sadyaḥ}
  \item \emph{Say it:} \emph{sadyaḥ}, once more
  \item \emph{Say it:} say \emph{sadyaḥ}
  \item \emph{Recall:} say the Sanskrit for a dream, then the Sanskrit for satisfaction, then say \emph{sadyaḥ} again
\end{itemize}

\begin{tcolorbox}[breakable,colback=teal!4,colframe=teal!35!black,title={Before you move on}]
What does \sk{सद्यः} mean? (\textquotedblleft{}At once\textquotedblright{}.) Say it once more.
\end{tcolorbox}
\section[\texorpdfstring{\sk{तत्क्षणम्} (tatkṣaṇam)}{तत्क्षणम् (tatkṣaṇam)}]{\sk{तत्क्षणम्} (tatkṣaṇam) — instantly}

\label{lesson:SA-C213-tatksanam}



\begin{quote}
Before the new one: say the Sanskrit for satisfaction, then the Sanskrit for at once.
\end{quote}

\subsection*{You'll want to know: \sk{तत्क्षणम्}}
\textbf{\sk{तत्क्षणम्}} — \emph{tatkṣaṇam} — \textquotedblleft{}instantly\textquotedblright{}.

\begin{grammarlens}[title={What you've built}]
One more word for this chapter.
\end{grammarlens}

\subsection*{Your turn}
\begin{itemize}
\raggedright
  \item \emph{Say it:} \emph{tatkṣaṇam}
  \item \emph{Say it:} \emph{tatkṣaṇam}, once more
  \item \emph{Say it:} say \emph{tatkṣaṇam}
  \item \emph{Recall:} say the Sanskrit for satisfaction, then the Sanskrit for at once, then say \emph{tatkṣaṇam} again
\end{itemize}

\begin{tcolorbox}[breakable,colback=teal!4,colframe=teal!35!black,title={Before you move on}]
What does \sk{तत्क्षणम्} mean? (\textquotedblleft{}Instantly\textquotedblright{}.) Say it once more.
\end{tcolorbox}
\section[\texorpdfstring{\sk{प्राक्} (prāk)}{प्राक् (prāk)}]{\sk{प्राक्} (prāk) — formerly, before}

\label{lesson:SA-C213-prak}



\begin{quote}
Before the new one: say the Sanskrit for at once, then the Sanskrit for instantly.
\end{quote}

\subsection*{You'll want to know: \sk{प्राक्}}
\textbf{\sk{प्राक्}} — \emph{prāk} — \textquotedblleft{}formerly, before\textquotedblright{}.

\begin{grammarlens}[title={What you've built}]
One more word for this chapter.
\end{grammarlens}

\subsection*{Your turn}
\begin{itemize}
\raggedright
  \item \emph{Say it:} \emph{prāk}
  \item \emph{Say it:} \emph{prāk}, once more
  \item \emph{Say it:} say \emph{prāk}
  \item \emph{Recall:} say the Sanskrit for at once, then the Sanskrit for instantly, then say \emph{prāk} again
\end{itemize}

\begin{tcolorbox}[breakable,colback=teal!4,colframe=teal!35!black,title={Before you move on}]
What does \sk{प्राक्} mean? (\textquotedblleft{}Formerly, before\textquotedblright{}.) Say it once more.
\end{tcolorbox}
\section[\texorpdfstring{\sk{अनन्तरम्} (anantaram)}{अनन्तरम् (anantaram)}]{\sk{अनन्तरम्} (anantaram) — afterwards, next}

\label{lesson:SA-C213-anantaram}



\begin{quote}
Before the new one: say the Sanskrit for instantly, then the Sanskrit for formerly.
\end{quote}

\subsection*{You'll want to know: \sk{अनन्तरम्}}
\textbf{\sk{अनन्तरम्}} — \emph{anantaram} — \textquotedblleft{}afterwards, next\textquotedblright{}.

\begin{grammarlens}[title={What you've built}]
One more word for this chapter.
\end{grammarlens}

\subsection*{Your turn}
\begin{itemize}
\raggedright
  \item \emph{Say it:} \emph{anantaram}
  \item \emph{Say it:} \emph{anantaram}, once more
  \item \emph{Say it:} say \emph{anantaram}
  \item \emph{Recall:} say the Sanskrit for instantly, then the Sanskrit for formerly, then say \emph{anantaram} again
\end{itemize}

\begin{tcolorbox}[breakable,colback=teal!4,colframe=teal!35!black,title={Before you move on}]
What does \sk{अनन्तरम्} mean? (\textquotedblleft{}Afterwards, next\textquotedblright{}.) Say it once more.
\end{tcolorbox}
\section[\texorpdfstring{\sk{यावत्} (yāvat)}{यावत् (yāvat)}]{\sk{यावत्} (yāvat) — as long as, until}

\label{lesson:SA-C213-yavat}



\begin{quote}
Before the new one: say the Sanskrit for formerly, then the Sanskrit for afterwards.
\end{quote}

\subsection*{You'll want to know: \sk{यावत्}}
\textbf{\sk{यावत्}} — \emph{yāvat} — \textquotedblleft{}as long as, until\textquotedblright{}.

\begin{grammarlens}[title={What you've built}]
One more word for this chapter.
\end{grammarlens}

\subsection*{Your turn}
\begin{itemize}
\raggedright
  \item \emph{Say it:} \emph{yāvat}
  \item \emph{Say it:} \emph{yāvat}, once more
  \item \emph{Say it:} say \emph{yāvat}
  \item \emph{Recall:} say the Sanskrit for formerly, then the Sanskrit for afterwards, then say \emph{yāvat} again
\end{itemize}

\begin{tcolorbox}[breakable,colback=teal!4,colframe=teal!35!black,title={Before you move on}]
What does \sk{यावत्} mean? (\textquotedblleft{}As long as, until\textquotedblright{}.) Say it once more.
\end{tcolorbox}
